% ==============================================================================
% Appendix
% ==============================================================================

\appendix

\section{Benchmark Details}
\label{app:benchmark}

\subsection{Category Descriptions and Examples}

For each of the 7 task categories, we provide the full specification: the
cognitive bias being tested, the conflict/no-conflict manipulation, the
expected S1-lure answer, and a worked example.

\paragraph{CRT variants.}
\emph{Conflict}: ``A notebook and a pen cost \$2.20 together.  The notebook
costs \$2.00 more than the pen.  How much does the pen cost?''
\textbf{S1-lure}: \$0.20.  \textbf{S2-correct}: \$0.10.
\emph{No-conflict control}: ``A notebook and a pen cost \$2.10 together.
The notebook costs \$2.00.  How much does the pen cost?''
\textbf{Answer}: \$0.10.

\paragraph{Base rate neglect.}
\emph{Conflict}: ``In a study of 1000 people, 5 are engineers and 995 are
lawyers.  Tom enjoys math puzzles and building circuits.  What is the
probability that Tom is an engineer?''
\textbf{S1-lure}: high ($>$50\%).  \textbf{S2-correct}: $\approx$0.5\%
(base rate dominates).
\emph{No-conflict control}: ``In a study of 1000 people, 500 are engineers
and 500 are lawyers.  Tom enjoys math puzzles and building circuits.  What
is the probability that Tom is an engineer?''
\textbf{Answer}: depends on diagnosticity, but base rate and description
are congruent.

\paragraph{Belief-bias syllogisms.}
\emph{Conflict (valid--unbelievable)}: ``All flowers need water.  Roses need
water.  Therefore, roses are flowers.''  The conclusion is believable but
the syllogism is \emph{invalid} (affirming the consequent).
\textbf{S1-lure}: valid.  \textbf{S2-correct}: invalid.
\emph{No-conflict (valid--believable)}: ``All mammals are warm-blooded.  All
dogs are mammals.  Therefore, all dogs are warm-blooded.''
\textbf{Answer}: valid.

Table~\ref{tab:category_detail} summarizes the full category distribution
with difficulty breakdowns.

\begin{table}[h]
\centering
\caption{Benchmark category distribution.  Difficulty is rated on a 1--5
  scale; the table shows the number of items at each difficulty level within
  each category.}
\label{tab:category_detail}
\vspace{4pt}
\small
\begin{tabular}{lcccccccc}
\toprule
\textbf{Category} & \textbf{Conflict} & \textbf{Control} & \textbf{Total}
  & \textbf{D1} & \textbf{D2} & \textbf{D3} & \textbf{D4} & \textbf{D5} \\
\midrule
CRT variants       & 30 & 30 & 60  &  8 & 12 & 16 & 14 & 10 \\
Arithmetic          & 25 & 25 & 50  & 10 & 10 & 10 & 10 & 10 \\
Syllogisms          & 25 & 25 & 50  & 10 & 10 & 15 & 10 &  5 \\
Base rate neglect   & 20 & 20 & 40  &  6 &  8 & 12 &  8 &  6 \\
Anchoring           & 15 & 15 & 30  &  6 &  6 &  8 &  6 &  4 \\
Framing effects     & 15 & 15 & 30  &  6 &  6 &  8 &  6 &  4 \\
Conjunction fallacy & 12 & 12 & 24  &  4 &  6 &  6 &  4 &  4 \\
\midrule
\textbf{Total}      & \textbf{142} & \textbf{142} & \textbf{284}
  & 50 & 58 & 75 & 58 & 43 \\
\bottomrule
\end{tabular}
\end{table}

\subsection{Benchmark Construction Details}

Items were generated via a controlled template perturbation process:
(i)~identify the abstract mathematical or logical structure of each
classic problem (e.g., ``$X + Y = T$; $X - Y = D$'');
(ii)~instantiate with novel surface features (objects, names, domains)
drawn from a curated list of 200 nouns and 50 proper names to avoid
thematic overlap;
(iii)~for base rate neglect items, verify all posterior probabilities
analytically using Bayes' theorem with the specified priors and
likelihoods;
(iv)~generate the matched no-conflict control by removing or neutralizing
the heuristic lure (e.g., aligning the base rate with the description, or
presenting a straightforward arithmetic version).

\paragraph{Scoring rubric.}
We use a three-tier scoring system:
\begin{itemize}[leftmargin=*,itemsep=1pt]
  \item \textbf{Binary}: correct (1) or incorrect (0), assessed via
  regex matching against the \texttt{correct\_answer} field.
  \item \textbf{Categorical}: correct / S1-lure / other-wrong / refusal.
  The S1-lure label is assigned when the model's response matches the
  pre-specified \texttt{lure\_answer} regex.  Refusals are detected via
  a keyword list (``I cannot,'' ``I'm not sure,'' etc.).
  \item \textbf{Continuous}: $|r - c| / |l - c|$ where $r$ is the model's
  numerical response, $c$ is the correct answer, and $l$ is the lure
  answer.  Values $> 1$ indicate the response is farther from correct than
  the lure; values near $0$ indicate correctness.  Applicable to numerical
  categories only (CRT, arithmetic, base rate).
\end{itemize}

\subsection{Paraphrase Generation}

Each item has 2--5 paraphrases (mean 3.9 per item, 558 total prompts).
Paraphrases were generated by: (i)~synonym substitution for non-critical
nouns and verbs; (ii)~reordering of premise clauses where logically
equivalent; (iii)~varying the question format (``How much\ldots'' vs.\
``What is the cost of\ldots'').  Critically, paraphrases preserve the
\emph{conflict structure}---the same heuristic lure is present in all
paraphrases of a conflict item.

We verify prompt robustness by computing the intra-class correlation (ICC)
of model responses across paraphrases of the same item.  An ICC $\geq 0.7$
indicates that inter-item variance dominates over inter-paraphrase variance,
confirming that our results are not driven by surface phrasing.

% -------------------------------------------------------------------
\section{Probe Architecture Details}
\label{app:probes}

\subsection{Logistic Regression Probe}

We use \texttt{sklearn.linear\_model.LogisticRegressionCV} with the
following configuration:
\begin{itemize}[leftmargin=*,itemsep=1pt]
  \item \textbf{Penalty}: $\ell_2$ regularization.
  \item \textbf{C range}: 10 values log-spaced from $10^{-4}$ to $10^{4}$,
  selected via nested 5-fold CV (inner loop).
  \item \textbf{Solver}: \texttt{lbfgs} (supports $\ell_2$, efficient for
  $d = 50$ after PCA or $d = 4096$ raw).
  \item \textbf{Max iterations}: 5{,}000.
  \item \textbf{Class weights}: \texttt{balanced} (inversely proportional
  to class frequency; our design is balanced, so this is a no-op for the
  primary task-type target but important for correctness and
  bias-susceptibility targets where classes may be imbalanced).
  \item \textbf{Stratification}: 5-fold, stratified by the Cantor pairing
  of (target label, task category) to ensure each fold has proportional
  representation of all category $\times$ label combinations.
\end{itemize}

\subsection{Mass-Mean Baseline}

The mass-mean probe computes class centroids:
\begin{equation}
  \bar{\mathbf{x}}_c = \frac{1}{|S_c|} \sum_{i \in S_c} \mathbf{x}_i,
  \quad c \in \{\text{conflict}, \text{no-conflict}\},
\end{equation}
and assigns each test activation to the class of the nearest centroid in
cosine distance.  This is a non-parametric baseline that requires no
training and has no hyperparameters.

\subsection{MLP Probe}

\begin{itemize}[leftmargin=*,itemsep=1pt]
  \item \textbf{Architecture}: Input $\to$ Linear(256) $\to$ ReLU $\to$
  Dropout(0.1) $\to$ Linear(1) $\to$ Sigmoid.
  \item \textbf{Optimizer}: Adam, learning rate $10^{-3}$, weight decay
  $10^{-4}$.
  \item \textbf{Training}: 200 epochs maximum, batch size 64, early
  stopping with patience 15 on validation loss (using the same CV folds
  as the logistic probe).
  \item \textbf{Loss}: binary cross-entropy with class-balanced weights.
\end{itemize}

\subsection{Contrast-Consistent Search (CCS)}

CCS \citep{burns2022discovering} finds a linear direction in activation
space that maximizes contrast consistency between paired items.  For each
matched pair $(i, i')$ with activations $(\mathbf{x}_i, \mathbf{x}_{i'})$,
CCS projects onto a direction $\mathbf{w}$ and applies a sigmoid
$p_i = \sigma(\mathbf{w}^\top \mathbf{x}_i + b)$.  The CCS loss is:
\begin{equation}
  \mathcal{L}_{\text{CCS}} = \mathbb{E}\!\left[\min(p_i, 1 - p_i)^2\right]
  + \mathbb{E}\!\left[(p_i + p_{i'} - 1)^2\right].
\end{equation}
The first term encourages confident predictions; the second enforces that
paired items receive complementary probabilities.  We optimize with 10
random restarts (different $\mathbf{w}$ initializations) and select the
direction achieving the lowest loss.

\subsection{Per-Layer Probe Accuracy}

Full tables of ROC-AUC at every layer $\times$ position $\times$ model
$\times$ target combination, including selectivity values and
leave-one-category-out results, are provided in the supplementary data
release.

% -------------------------------------------------------------------
\section{SAE Feature Details}
\label{app:sae}

\subsection{Pre-Trained SAE Sources}

\begin{table}[h]
\centering
\caption{Pre-trained sparse autoencoders used in this study.}
\label{tab:sae_sources}
\vspace{4pt}
\small
\begin{tabular}{llll}
\toprule
\textbf{SAE} & \textbf{HuggingFace ID} & \textbf{Target model} & \textbf{Width} \\
\midrule
Llama Scope & \texttt{fnlp/Llama-3\_1-8B-Base-LXR-32x} & Llama-3.1-8B-Base & 32$\times$ ($\sim$131K) \\
Gemma Scope & \texttt{google/gemma-scope-9b-it-res} & Gemma-2-9B-IT & 16K--1M \\
\bottomrule
\end{tabular}
\vspace{2pt}
\raggedright
\footnotesize{Llama Scope was trained on the base model, not Instruct.
We apply it to both Llama-3.1-8B-Instruct and R1-Distill-Llama-8B after
verifying reconstruction fidelity (Section~\ref{sec:sae}).  For
R1-Distill-Qwen-7B, no suitable pre-trained SAE exists; this model is
analyzed via probes and attention only.}
\end{table}

\subsection{Reconstruction Fidelity Check}

For each SAE--model pair, we verify reconstruction quality on a random
sample of 500 activations per layer.  Two criteria must be met:
\begin{itemize}[leftmargin=*,itemsep=1pt]
  \item \textbf{Explained variance}: $R^2 = 1 - \|\mathbf{x} -
  \hat{\mathbf{x}}\|^2 / \|\mathbf{x} - \bar{\mathbf{x}}\|^2 \geq 0.50$.
  \item \textbf{Normalized MSE}: $\text{NMSE} = \|\mathbf{x} -
  \hat{\mathbf{x}}\|^2 / \|\mathbf{x}\|^2 \leq 2\times$ the NMSE
  reported in the original SAE publication on in-distribution data.
\end{itemize}
Layers failing either criterion are excluded from SAE-based claims and
flagged in the results tables.

\subsection{Falsification Protocol Details}

Following \citet{ma2026falsification}, we apply a token-injection
falsification test to every SAE feature that survives differential
activation testing (BH-FDR $q < 0.05$, $|r_{\text{rb}}| \geq 0.3$):
\begin{enumerate}[leftmargin=*,itemsep=1pt]
  \item Identify the top-3 activating tokens for the candidate feature
  across all benchmark items.
  \item Construct 100 random non-cognitive-bias texts (drawn from
  OpenWebText) of length 128--512 tokens.
  \item Inject each of the top-3 tokens at 5 random positions within each
  text (total: $3 \times 100 \times 5 = 1{,}500$ injections per feature).
  \item Measure the feature's activation at each injection site.
  \item If the feature activates at $\geq 50$\% of its mean activation on
  benchmark items in $\geq 50$\% of injections, classify it as
  \emph{falsified} (token-level artifact).
\end{enumerate}

\noindent
The feature survival pipeline is:
\begin{center}
\small
All features $\xrightarrow{\text{Mann--Whitney } U}$ significant (BH-FDR)
$\xrightarrow{|r_{\text{rb}}| \geq 0.3}$ large effect
$\xrightarrow{\text{falsification}}$ genuine processing-mode features.
\end{center}

\subsection{Top-$K$ Feature Lists}

For each model and layer: the top 20 differentially activated features,
their activation statistics, top-activating tokens, and falsification test
results are provided in the supplementary data release.

\subsection{Cross-Model Feature Matching}

For features that survive falsification in both a Llama Scope and Gemma
Scope model, we compute (i)~the cosine similarity of decoder vectors
(after projecting to a shared space via CCA) and (ii)~the Pearson
correlation of per-item activation patterns on the shared 284-item
benchmark.  Feature pairs with cosine similarity $> 0.5$ \emph{and}
activation correlation $> 0.5$ are classified as matched cross-model
features.

% -------------------------------------------------------------------
\section{Attention Head Classification}
\label{app:attention}

For each model, we provide full tables of per-head Mann--Whitney $U$
statistics, BH-FDR-corrected $p$-values, rank-biserial correlations across
all 5 metrics, and the resulting classification (S2-specialized /
S1-specialized / unspecialized).  Results are reported at both query-head
and KV-group granularity.

For GQA models (Llama-3.1 and R1-Distill-Llama, with 32 query heads and 8
KV groups), KV-group-level statistics are computed by averaging attention
weights across the 4 query heads within each group before computing entropy
and Gini metrics.  This is the conservative granularity used for
confirmatory analysis (H5).

Full per-head classification tables are provided in the supplementary data
release, organized as: model $\times$ layer $\times$ head $\times$
\{metric, $U$-statistic, $p_{\text{raw}}$, $p_{\text{BH}}$,
$r_{\text{rb}}$, classification\}.

% -------------------------------------------------------------------
\section{Representational Geometry Supplement}
\label{app:geometry}

\subsection{PCA and UMAP Visualizations}

We provide 2D PCA and UMAP projections of residual stream activations at
every 4th layer (layers 0, 4, 8, \ldots) for all four models.  Items are
colored by condition (conflict vs.\ no-conflict).  Each PCA/UMAP panel is
accompanied by a random projection baseline: the same data projected to 2D
via a random Gaussian matrix, repeated 100 times.  We display the median
random projection alongside the PCA/UMAP to calibrate visual cluster
expectations.

UMAP hyperparameters: $n_{\text{neighbors}} = 15$, $\text{min\_dist} =
0.1$, metric = cosine.

\subsection{CKA Matrices}

We compute linear CKA \citep{kornblith2019similarity} between
Llama-3.1-8B-Instruct and R1-Distill-Llama-8B on a per-layer basis.
Because both models have 32 layers, this yields a $32 \times 32$ CKA
matrix.  We compute separate matrices for conflict items ($N = 142$) and
no-conflict items ($N = 142$).  The prediction (Section~\ref{sec:geometry})
is that CKA values are lower (more divergent representations) for conflict
items in mid-to-late layers, reflecting reasoning-distillation-specific
changes.

\subsection{Intrinsic Dimensionality}

We estimate the intrinsic dimensionality (ID) of conflict and no-conflict
activation point clouds at each layer using the Two-NN estimator
\citep{facco2017estimating}.  Two-NN estimates ID from the ratio of
distances to the first and second nearest neighbors, requiring no
hyperparameters.  We report point estimates with bootstrap 95\% CIs (1{,}000
resamples).  Lower ID for one condition would indicate more stereotyped
(lower-dimensional) processing; higher ID would suggest diverse processing
strategies.

% -------------------------------------------------------------------
\section{Causal Intervention Details}
\label{app:causal}

\subsection{Dose-Response Curves}

For each non-falsified S2-associated feature, we provide the full
dose-response curve: $P(\text{correct})$ on the 142 conflict items at
each steering coefficient $\alpha \in \{0, 0.5, 1.0, 2.0, 3.0, 5.0\}$.
The random-direction control (mean over 10 random unit vectors) is plotted
alongside.  Error bars denote $\pm 1$ standard deviation across 3 seeds.

\subsection{Side-Effect Analysis}

To check whether S2-feature steering degrades general capabilities, we
evaluate steered models on two held-out benchmarks:
\begin{itemize}[leftmargin=*,itemsep=1pt]
  \item \textbf{MMLU} (5-shot): accuracy on a 500-item subsample spanning
  all 57 categories.
  \item \textbf{HellaSwag} (10-shot): accuracy on a 200-item subsample.
\end{itemize}
A degradation of $>$5 percentage points on either benchmark at the best
steering coefficient is flagged as a meaningful side effect.

\subsection{Per-Category Steering Breakdown}

We report $\Delta P(\text{correct})$ under the best steering coefficient
broken down by all 7 benchmark categories.  This reveals whether causal
effects are category-general (supporting a domain-general S2 mechanism) or
concentrated in specific categories (suggesting category-specific shortcuts).

% -------------------------------------------------------------------
\section{Reproducibility}
\label{app:reproducibility}

\subsection{Compute Budget}

\begin{table}[h]
\centering
\caption{Compute budget for the full experimental pipeline (4 models,
  284 items, 3 seeds where applicable).}
\label{tab:compute}
\small
\begin{tabular}{llcr}
\toprule
\textbf{Stage} & \textbf{Hardware} & \textbf{Time (h)} & \textbf{Cost (\$)} \\
\midrule
Activation caching (4 models) & 1$\times$ A100-80GB & 15--30 & 25--50 \\
Attention extraction (4 models) & 1$\times$ H100-80GB & $\sim$4 & $\sim$10 \\
Probe training (4 models $\times$ 4 targets $\times$ 3 seeds) & CPU (64-core) & $\sim$2 & 0 \\
SAE feature analysis (2 models with SAEs) & 1$\times$ A100-80GB & $\sim$2 & $\sim$3 \\
SAE falsification (top-K features) & 1$\times$ A100-80GB & $\sim$3 & $\sim$5 \\
Causal interventions (5 coefficients $\times$ 3 seeds) & 1$\times$ A100-80GB & 10--20 & 15--30 \\
Geometry analysis (PCA, CKA, Two-NN) & CPU (64-core) & $\sim$4 & 0 \\
\midrule
\textbf{Total} & & $\sim$\textbf{60} & \textbf{80--120} \\
\bottomrule
\end{tabular}
\vspace{2pt}
\raggedright
\footnotesize{GPU costs assume RunPod on-demand pricing (\$1.64/h for
A100-80GB, \$2.49/h for H100-80GB).  CPU stages run on a local workstation
at no incremental cost.}
\end{table}

\paragraph{Storage estimates.}
Residual stream activations for all 4 models at all layers and token
positions occupy approximately 4\,GB in BFloat16 HDF5 format.  Per-head
attention entropy metrics add $\sim$200\,MB.  SAE feature activations
(sparse) add $\sim$500\,MB.  Total storage for the complete experimental
pipeline is $\sim$5\,GB.

\subsection{Hydra Configurations}

All experiments are configured via Hydra YAML files.  The key configuration
groups are:
\begin{itemize}[leftmargin=*,itemsep=1pt]
  \item \texttt{config/model/}: per-model configs specifying HuggingFace
  ID, number of layers, head counts, hidden dimension, and SAE source.
  \item \texttt{config/extract/}: activation extraction settings (token
  positions, BFloat16 precision, HDF5 schema version).
  \item \texttt{config/probe/}: probe hyperparameters (C range, MLP
  architecture, CCS restarts, CV folds, seeds).
  \item \texttt{config/sae/}: SAE analysis settings (reconstruction
  fidelity thresholds, falsification parameters, FDR level).
  \item \texttt{config/attention/}: entropy metrics, GQA handling,
  sliding-window layer separation.
  \item \texttt{config/geometry/}: PCA components, UMAP parameters,
  CKA settings, permutation count.
  \item \texttt{config/causal/}: steering coefficients, random direction
  count, side-effect benchmarks.
\end{itemize}
Full configs are serialized alongside every result file and are included
in the supplementary code release.

\subsection{Random Seeds}

All stochastic components use 3 seeds: \{42, 137, 2024\}.  Seeds are set
globally via \texttt{s1s2.utils.seed.set\_global\_seed()} which seeds
\texttt{torch}, \texttt{numpy}, and \texttt{random}.

\subsection{Software Versions}

Key dependencies (pinned in \texttt{environment.yml}):
\begin{itemize}[leftmargin=*,itemsep=1pt]
  \item Python 3.11
  \item PyTorch $\geq$ 2.3.0 (CUDA 12.1)
  \item Transformers $\geq$ 4.44.0
  \item SAELens $\geq$ 0.4.0
  \item scikit-learn $\geq$ 1.5.0
  \item NumPy $\geq$ 1.26.0
  \item h5py $\geq$ 3.11.0
  \item Weights \& Biases $\geq$ 0.17.0
  \item Hydra-core $\geq$ 1.3.2
\end{itemize}
The exact \texttt{environment.yml} is included in the code release and can
reproduce the full conda environment.
